%% LyX 2.4.4 created this file.  For more info, see https://www.lyx.org/.
%% Do not edit unless you really know what you are doing.
\documentclass[english]{article}
\usepackage[T1]{fontenc}
\usepackage[latin9]{inputenc}
\usepackage{units}
\usepackage{enumitem}
\usepackage{amsmath}
\usepackage{amsthm}
\usepackage{cancel}
\usepackage{geometry}
\geometry{verbose,tmargin=1in,bmargin=1in,lmargin=1in,rmargin=1in}

\makeatletter
%%%%%%%%%%%%%%%%%%%%%%%%%%%%%% Textclass specific LaTeX commands.
\numberwithin{equation}{section}
\numberwithin{figure}{section}
\newlength{\lyxlabelwidth}      % auxiliary length 

%%%%%%%%%%%%%%%%%%%%%%%%%%%%%% User specified LaTeX commands.
\usepackage{fancyhdr}
\usepackage{lastpage}
\pagestyle{fancy}
\usepackage{enumitem}
\usepackage{ifthen}
%sets math fonts to be more readable
\everymath=\expandafter{\the\everymath\displaystyle}
%\DeclareMathSizes{10}{10}{10}{10}
\usepackage{multicol}

\usepackage{tikz}
\usetikzlibrary{plotmarks}
\usepackage{pgfplots}
\pgfplotsset{compat=newest}

\usepackage{calc}
\usepackage[nomessages]{fp}

\usepackage{advdate}    % Advancing/saving dates
\usepackage{datetime}   % Dates formatting
\usepackage{datenumber} % Counters for dates
% Added by lyx2lyx
\setlength{\parskip}{\smallskipamount}
\setlength{\parindent}{0pt}

\makeatother

\usepackage{babel}
\begin{document}
\renewcommand{\author}{Dr.\,James Ashe}

\newcommand{\classNum}{137}

\newcommand{\classSec}{B}

\newcommand{\nameType}{Name:}

\newcommand{\examType}{worksheet 4.1-4.3}

\renewcommand{\date}{\formatdate{22}{4}{2013}}

%%First page's header%%%%%%%%%%%%%%%%%%%%%%
\fancypagestyle{firststyle} %%header settings for first pagr
{
	\fancyhead[L]{MTH \classNum{}(\classSec{}) \author{}\\
	\textbf{\examType{}} \date{}}
	\fancyhead[C]{\textbf{\nameType}}
	\setlength{\headsep}{14pt}
}
%%%%%%%%%%%%%%%%%%%%%%%%%%%%%%%%%%%%%%%%%%

%%Next page's header%%%%%%%%%%%%%%%%%%%%%%%
\setlist{nolistsep}
\lhead{\classNum{}(\classSec{})} 
\chead{\textbf{\nameType}} 
\cfoot{\thepage{}~of~\pageref{LastPage}} 
\setlength{\headsep}{5pt} 
%%%%%%%%%%%%%%%%%%%%%%%%%%%%%%%%%%%%%%%%%%%

%%Various settings; see the preamble for other settings%%%%%
%%\setenumerate[0]{leftmargin=12pt,itemindent=0pt,labelwidth=0pt}
\renewcommand{\labelenumi}{\textbf{\arabic{enumi}.}}
\renewcommand{\labelenumii}{\textbf{\alph{enumii}.}}
\thispagestyle{firststyle}
%%%%%%%%%%%%%%%%%%%%%%%%%%%%%%%%%%%%%%%%%%
\newcommand{\xmax}{10}
\newcommand{\xmin}{-10}
\newcommand{\xstep}{0} %must be xmin+n
\newcommand{\ymax}{10}
\newcommand{\ymin}{-10}
\newcommand{\ystep}{0} %must be ymin+n

\newcommand{\Dspace}{\vspace{1cm}}
\begin{enumerate}
\item Evaluate the following:

\begin{enumerate}
\item $-2^{-3}$\Dspace
\item $\left(-\frac{2}{3}\right)^{-2}$\Dspace
\item $8^{-\nicefrac{4}{3}}$\Dspace
\item $f(x)=5^{2-3x}$, find $f(1)$.\Dspace
\end{enumerate}
\end{enumerate}
\begin{enumerate}[resume]
\item Consider the polynomial $f(x)$ shown in the graph.\\
\renewcommand{\ymax}{125}
\renewcommand{\ymin}{-250}
%%%%%%%
\begin{tikzpicture} 
\begin{axis}[height=5cm, 
	width=5cm, 
	axis lines=middle,
	minor x tick num={4},
	minor y tick num={4},
	%xtick={0,50,...,250},
	%ytick={0,10,20,...,40},
	%grid=both,
	%axis line style={-latex}, 
	xmin=\xmin,xmax=\xmax, 
	ymin=\ymin,ymax=\ymax,
	%restrict y to domain=-1:10,
	%no markers, 
	xlabel={$x$},ylabel={$y$}, 
	%every axis x label/.append style={anchor=north}, 
	%every axis y label/.append style={anchor=east}
	]
	
%%\addplot[color=red,solid,mark=*] coordinates {(0,2)};
	\addplot[domain=\xmin+2:\xmax-1.5,thick,samples=1000,draw=red,<->]{(x+7)*(x-8)*(x+1)} node[left] {$f(x)$};
	
\end{axis}
\end{tikzpicture}

\begin{enumerate}
\item What are the zeros of $f(x)$?\Dspace\Dspace
\item Does $f(x)=(x-7)(x+8)(x+1)$?\Dspace
\end{enumerate}
\end{enumerate}
\begin{itemize}
\item If the factor $x-c$ occurs $k$ times in the complete factorization
of a polynomial $P(x)$, then $c$ is called a root/zero of $P(x)$
with multiplicity $k$.
\item If $P(x)$ is polynomial of degree $n$ then it has at most $n$ real
roots/zeros (counting multiplicity).
\item \textbf{Descartes's Rule of Signs: }If $P(x)$ is a polynomial written
in descending order then,

\begin{itemize}
\item The number of $P(x)$'s \emph{positive} roots is equal to the number
of $P(x)$'s sign changes or less than that by an even number.
\item The number of $P(x)$'s \emph{negative} roots is equal to the number
of $P(-x)$'s sign changes or less than that by an even number.\Dspace
\end{itemize}
\end{itemize}
\begin{enumerate}[resume]
\item Let $P(x)=x^{4}+2x^{3}-5x^{2}-6x$.

\begin{enumerate}[resume]
\item What is the largest possible number of roots for $P(x)$?\Dspace\Dspace
\item What are the possible number of positive roots for$P(x)$?\Dspace\Dspace
\item What are the possible number of negative roots for $P(x)$?\Dspace\Dspace
\end{enumerate}
\end{enumerate}

\end{document}
